% Ellipse de la hessienne lissée en un pixel d'une crête sombre.
% On trace le niveau {u : u^T |H_sigma| u = 1} avec |H| = |l1| v1 v1^T + |l2| v2 v2^T :
% c'est bien une ellipse quels que soient les SIGNES des valeurs propres (le niveau
% de H lui-même serait une hyperbole dès que l1 < 0, ce qui arrive sur une crête où
% l1 ~ 0). Demi-axes 1/sqrt(|l_k|) portés par les vecteurs propres v_k.
% Crête : |l1| petit (grand axe, longitudinal, angle theta) et l2 > 0 grand (petit axe).
    \begin{tikzpicture}[x=1cm,y=1cm,scale=1.0,every node/.style={scale=1.0},
        crete/.style={line width=8pt, gris_fonce_inria!55, line cap=round},
        vecp/.style={-{Latex[length=2.1mm]}, line width=1pt},
        filet/.style={line width=0.8pt}]

      \begin{scope}[rotate=27]
        \draw[crete] (-2.15,0)--(2.15,0);
        \draw[inria-rouge, line width=0.9pt, fill=inria-rouge!7]
              (0,0) ellipse (1.62 and 0.46);
        % longueurs corrigées de la pointe de flèche : les pointes tombent en 1,62 et 0,46
        \draw[vecp, inria-rouge]      (0,0)--(1.6925,0);
        \draw[vecp, inria-rouge!62]   (0,0)--(0,0.5325);
        \draw[inria-rouge!45, dashed, line width=0.5pt] (1.62,0)--(2.24,0);
        \node[font=\scriptsize, inria-rouge, fill=white, inner sep=0.8pt] at (0.80,0) {$v_1$};
        \node[font=\scriptsize, inria-rouge!62, right=1pt] at (0.02,0.30) {$v_2$};
      \end{scope}

      % axe de référence et angle theta, tracés HORS de l'ellipse
      \draw[gray!65, dashed, line width=0.5pt] (0,0)--(2.10,0);
      \draw[gris_fonce_inria!75, line width=0.6pt] (1.90,0) arc (0:27:1.90);
      \node[font=\scriptsize] at (1.72,0.44) {$\theta$};
      \node[circle,fill=inria-2024-bleu-canard,inner sep=0pt,minimum size=3.6pt] (px) at (0,0) {};
      \node[font=\scriptsize, below left=0pt and 1pt of px] {$x$};

      % demi-axes, rattachés par des filets qui sortent franchement de la bande
      \draw[inria-rouge, filet] (1.4434,0.7355)--(1.62,1.06);
      \node[font=\scriptsize, inria-rouge, anchor=south west] at (1.56,1.05)
            {$\dfrac{1}{\sqrt{|\lambda_1|}}$};
      \draw[inria-rouge!62, filet] (-0.2088,0.4099)--(-0.72,0.86);
      \node[font=\scriptsize, inria-rouge!62, anchor=south east] at (-0.68,0.83)
            {$\dfrac{1}{\sqrt{|\lambda_2|}}$};
      \draw[gris_fonce_inria!70, filet] (-1.72,-0.92)--(-1.86,-1.22);
      \node[font=\scriptsize, text=gris_fonce_inria, anchor=north, align=center] at (-1.86,-1.22)
            {crête\\sombre};
    \end{tikzpicture}
